\section{工业 4.0 中的控制技术}

\subsection{控制技术方式}

在传统工厂中，机器主要依赖固定硬件和本地控制。例如一台 CNC 机床、一台机器人、一套 PLC 控制柜，各自完成自己的任务。而工业 4.0 的核心趋势是把这些设备联网、数字化、软件化，使它们能够交换数据、被监控、被优化，甚至在一定程度上自主调整。

我们用 Software Defined Manufacturing (SDM) 来描述这个趋势：工业自动化正在从 hardware-oriented 转向 software-oriented。也就是说，以前重点是专用硬件和固定功能；未来重点是软件模块、虚拟化、云技术、数字孪生、AI 和自适应系统。

\subsection{工业控制系统的类型}

我们把常见控制系统分成 NC、RC、MC、SPS 和 MES。

工厂不是一个单一控制器控制所有东西，而是按层级组织的。越靠近现场设备，越强调实时性；越靠近管理层，越强调计划、调度、数据和业务流程。控制层级从 E1 到 E7：
\begin{itemize}
    \item E1 是过程层，例如传感器和执行器
    \item E2 是单机控制
    \item E3 是机器控制
    \item E4 是单元/工作站层
    \item E5 是生产线或车间管理层
    \item E6 是计划层
    \item E7 是工厂控制层
\end{itemize}

越往下，越接近物理过程，实时性要求越高。

NC 是 Numerical Control，数控，典型对象是 CNC 机床。它的任务是根据数字程序控制刀具或工件的运动。例如铣削、车削、钻孔时，程序给出几何路径、进给速度、主轴转速等，NC 将这些信息转换成各个轴的运动命令。NC 的特点包括：
\begin{itemize}
    \item 通过数字进行控制
    \item 按 DIN 66025 编程
    \item 处理多个自由度
    \item 将笛卡尔坐标转换为轴坐标
\end{itemize}

NC 常用于工具机、模具加工和机械制造，并追求高绝对精度。

RC 是 Robot Control，机器人控制。它控制的是工业机器人，特别是串联机械臂。机器人控制的重点不是切削路径，而是末端执行器或夹爪的空间运动。机器人程序可以通过离线编程、Teach-In 或 Playback 得到。它需要在关节坐标和空间坐标之间做坐标变换。

RC 常用于汽车制造、装配和搬运，重点通常是高重复精度，而不一定是像 CNC 那样的高绝对精度。

MC 是 Motion Control，运动控制。它关注多个轴之间的同步运动，例如印刷机、包装机、飞锯、输送带同步、电子凸轮等。它的核心不是“走一条复杂加工轨迹”，而是让多个驱动轴按照严格时序协同运动。比如一条高速包装线中，输送带、切刀、推料机构必须在毫秒级时间内协调，否则产品会错位。

SPS 是 Speicherprogrammierbare Steuerung，也就是 PLC。它是自动化中最基础、最常见的控制器。PLC 读取传感器输入，执行逻辑程序，然后控制阀门、电机、继电器、气缸等执行器。

SPS 按 DIN EN 61131-3 编程，处理模拟和数字信号，采用循环程序执行，直接连接传感器和执行器，并且结构紧凑、鲁棒。它广泛用于制造技术、过程工业和建筑自动化。

MES 是 Manufacturing Execution System，制造执行系统。它的位置高于机器控制，低于企业资源计划系统。MES 不直接以毫秒级控制电机，而是管理生产执行过程。例如订单执行、工位占用、生产状态、物料流、仓储、质量保证、数据采集、追溯和序列化。可以把 MES 理解为“车间层面的操作系统”：它知道哪个产品在什么工位、完成了什么步骤、用了什么材料、是否合格、下一步该去哪里。

图 \ref{fig:steuerungebene} 展示了这7个层级以及对应的功能。

\begin{figure}[htbp]
  \centering
  \includegraphics[width=0.8\textwidth]{steuerungebene.png}
  \caption{工厂的控制层级}
  \label{fig:steuerungebene}
\end{figure}

\subsection{控制架构}

我们不能只讲“有哪些控制器”，还要讲控制功能放在哪里。一个自动化系统通常有人机界面 HMI、逻辑控制和运动控制三类软件功能。我们用 OMAC 分类，把控制架构分成 PC-based、controller-based、drive-based 和 cloud-based。
\begin{itemize}
    \item PC-based 架构把 HMI、逻辑和运动控制集中在一台工业 PC 上。优点是软件集中、计算能力强、易于扩展和集成；缺点是实时性、可靠性和系统隔离必须通过合适的实时机制来保证。
    \item Controller-based 架构把逻辑和运动控制放在专用控制器中，HMI 则可能运行在单独的面板或 PC 上。这是工业中很常见、很稳健的架构。专用控制器更接近传统 PLC/CNC 控制方式，实时性和可靠性通常更容易保证。
    \item Drive-based 架构把一部分逻辑和运动控制下放到驱动器中。例如某些智能伺服驱动器不仅执行电流环、速度环，还能执行局部运动逻辑。这样可以减少上位控制器负担，提高局部响应速度。
    \item Cloud-based 架构则把部分功能放到边缘云或云端。它适合监控、优化、数据分析、预测性维护等任务，但不适合把严格毫秒级闭环控制完全放到远端云端，因为网络延迟和确定性难以保证。
\end{itemize}

\subsection{NC内部的信息流}

数控系统内部如何从加工任务变成实际轴运动。加工程序首先描述几何轮廓和工艺参数；控制系统解析程序；然后进行轨迹规划、插补、坐标变换；再生成每个轴的设定值；最后驱动器和伺服系统控制电机运动。这里的关键概念是“从抽象几何任务到实时轴控制”的转换。比如程序里写的是刀具从 A 点走到 B 点，但机器真正执行时，需要每个采样周期给 X、Y、Z 轴分别发送位置或速度指令。

\subsection{工业通信技术}
\label{subsec:industrialcommunication}

工业 4.0 不只是机器能动起来，还要求机器之间、机器与上位系统之间能够可靠交换数据。因此就得用到 ISO/OSI 七层模型，并把它和工业通信联系起来。

七层包括：
\begin{enumerate}
    \item 物理层
    \item 数据链路层
    \item 网络层
    \item 传输层
    \item 会话层
    \item 表示层
    \item 应用层
\end{enumerate}

例如：Ethernet、TSN 属于较底层的数据链路相关技术；IP 位于网络层；TCP/UDP 位于传输层；OPC-UA 属于应用层。

OPC-UA 可以理解为工业系统之间交换数据的标准语言。它解决的问题是：不同厂家、不同控制器、不同软件系统之间如何以统一方式描述和访问数据。例如一个控制器中的轴位置、温度、报警状态、设备 ID，如果没有统一语义，每个厂家都用自己的格式，系统集成会非常困难。OPC-UA 提供统一的信息模型和通信机制，使设备数据可以被上位系统、MES、SCADA 或其他应用读取和订阅。

TSN 是 Time Sensitive Networking，时间敏感网络。它的目标是让普通以太网具备确定性通信能力。普通以太网适合传文件、网页、邮件，但不能保证某个控制数据包一定在严格时间窗口内到达。工业控制尤其是运动控制需要确定性延迟、低抖动和同步。TSN 通过时间同步、优先级、带宽预留和时间窗口等机制，使工业网络更适合实时控制。课件用 OPC-UA PubSub over TSN 举例说明：应用层通过 OPC-UA 发布/订阅数据，底层通过 TSN 保证网络传输的时间确定性。

\subsection{实时操作系统}

实时系统的重点不是“速度快”，而是“在规定时间内可靠完成”。一个控制任务如果必须每 1 ms 更新一次轴位置，那么它不是平均 0.5 ms 完成就够了，而是每一次都必须在 deadline 之前完成。课件的学习检查问题中特别提到：进程调度状态、preemptive 与 non-preemptive scheduling、实时的含义，以及中断处理的时间行为。

Preemptive scheduling 是抢占式调度。高优先级任务到来时，可以打断低优先级任务。这适合实时控制，因为紧急任务可以及时执行。Non-preemptive scheduling 是非抢占式调度，一个任务开始后通常要运行到主动让出 CPU，这种方式简单，但对严格实时任务不利。中断处理也很重要，因为传感器、通信接口、驱动器反馈等事件可能需要立即响应。实时操作系统要保证中断延迟、任务切换延迟和周期任务执行时间可预测。